\chapter{Frontend, Names, and Modules}

The front end turns source text into a source-located AST and then into a
DefId-resolved IR. This part of the implementation is non-trivial because
\einlang{} is not only an expression language: it has modules, imports,
function hoisting, index scopes, named rest patterns, and local derivative
syntax.

\section{Parsing and AST Construction}

The parser in \pathcode{src/einlang/frontend/parser.py} wraps a shared Lark LALR
parser loaded from \pathcode{src/einlang/frontend/grammar.lark}. The shared
parser is keyed by cache path so repeated compilations do not rebuild grammar
tables. Disk parser caches are disabled by default, but the code can use an
explicit cache path when configured. The parser propagates source positions and
converts Lark errors into \einlang{} parse errors with source locations.

AST construction is handled by frontend transformers under
\pathcode{src/einlang/frontend/transformers}. The transformer layer is split by
concern: base expression and statement construction, function-specific logic,
literals, and string interpolation. This split matters because later compiler
passes assume that syntax has already been normalized into the AST node classes
in \pathcode{src/einlang/shared/nodes.py}.

The transformer is not just a mechanical parse-tree adapter. It constructs
source spans, normalizes equivalent surface forms, attaches structured nodes for
patterns and type annotations, and records interpolation pieces so the backend
can later evaluate a string expression without reparsing source text. Lark
therefore supplies the parsing engine, but \einlang{} still owns the language
grammar and the AST shape consumed by later compiler phases.

\section{Grammar Coverage}

The grammar contains the executed tensor subset and the broader language shape.
It parses:

\begin{itemize}
\item declarations: \texttt{fn}, \texttt{pub fn}, \texttt{@fn},
\texttt{let}, \texttt{const}, and indexed Einstein declarations;
\item expressions: arithmetic, logic, ranges, casts, calls, member access,
array access, array literals, comprehensions, tuples, blocks, \texttt{if},
\texttt{match}, lambdas, pipelines, and \texttt{try};
\item tensor syntax: rectangular access, reductions, explicit index ranges,
literal indices, and named rest indices such as \texttt{..batch};
\item modules: \texttt{use}, wildcard imports, list imports, aliases,
\texttt{mod}, public modules, and inline modules;
\item types: primitive, rectangular, jagged, tuple, function, generic, and
named types.
\end{itemize}

The important implementation choice is that tensor notation is not parsed as a
separate mini-language. It participates in the same expression grammar as the
rest of the source program.

\section{DefId-Based Identity}

Names become semantic identifiers during name resolution. The DefId system in
\pathcode{src/einlang/shared/defid.py} follows the rustc pattern: a definition
identifier is a pair \texttt{(krate, index)}. User code uses crate 0, and
builtins use a separate builtin crate with fixed DefIds for values such as
\texttt{print}, \texttt{assert}, \texttt{len}, \texttt{shape}, \texttt{sum},
\texttt{max}, and \texttt{min}.

This design removes ambiguity from later stages. Once name resolution has run,
runtime lookup and compiler analyses use DefIds rather than strings. Strings
remain for diagnostics and debugging, but they do not decide identity. This is
especially important for index variables and local variables because the same
spelling may appear in different scopes.

\section{Name Resolution and Scope}

The name-resolution pass in \pathcode{src/einlang/passes/name_resolution.py}
allocates DefIds, resolves use sites, handles builtins, assigns DefIds to
parameters and locals, and wires functions and modules into the compilation
context. Function differential rules are first merged into ordinary functions by
\pathcode{src/einlang/passes/merge_diff_rules.py}, so the compiler sees one
function definition with attached differential behavior instead of two
unrelated definitions.

The scope system in \pathcode{src/einlang/shared/scope.py} tracks scope kind and
binding type. Local variable DefIds are allocated without inserting them into
the global symbol table; item-like definitions such as functions, constants,
modules, and builtins are registered in the resolver tables. That separation is
what lets local shadowing work without corrupting module-level lookup.

\subsection{Resolution Algorithm}

The resolver is a two-level identity algorithm. Item-like names are installed in
module-level maps, while local names are installed only in lexical scopes. Every
use site is rewritten from a string name to a DefId before later passes run.
The key invariant is simple: after resolution, semantic equality is DefId
equality.

\begin{lstlisting}[style=pythoncode,caption={DefId-based name resolution.}]
resolve_program(ast, tcx):
    group_consecutive_einstein_declarations(ast)
    install_fixed_builtins(tcx.resolver)
    load_and_link_imported_modules(ast, tcx.module_loader, tcx.symbol_linker)

    # Hoist item identities before resolving bodies.
    for item in ast.top_level_items:
        if item is function or constant or module:
            did = resolver.allocate_for_item(item.module_path, item.name)
            item.defid = did
            resolver.register(did, item)
            current_module.define(item.name, Binding(did, ITEM))

    for item in ast.top_level_items:
        resolve_item_body(item)

resolve_item_body(item):
    with scope(kind=FUNCTION_OR_MODULE):
        for parameter in item.parameters:
            did = resolver.allocate_for_local()
            parameter.defid = did
            scope.define(parameter.name, Binding(did, PARAMETER))
        resolve_expr(item.body)

resolve_expr(expr):
    case local binding:
        resolve_expr(expr.value)
        did = resolver.allocate_for_local()
        expr.defid = did
        scope.define(expr.name, Binding(did, LOCAL))
    case identifier:
        binding = scope.lookup(expr.name) or module_lookup(expr.name)
        expr.defid = binding.defid
    case indexed clause:
        with scope(kind=INDEX):
            allocate DefIds for output indices and where-bindings
            resolve body and guard expressions in that scope
    case nested block:
        with scope(kind=BLOCK):
            resolve statements in source order
\end{lstlisting}

This algorithm explains why shadowing is safe. A local \texttt{i} in one
Einstein clause and another local \texttt{i} in a different clause may share a
printed name, but their DefIds differ. Range, shape, type, lowering, and
runtime lookup therefore do not need to guess which \texttt{i} a later node
means.

The listing is longer than a conventional symbol-table sketch because the
resolver is also the first pass that makes tensor syntax semantically precise.
Top-level hoisting gives functions and constants stable identities before their
bodies are visited, which permits forward references and module-level cycles to
be diagnosed at the right boundary. Local allocation is intentionally separate:
parameters, block locals, and index variables live in lexical scopes and do not
pollute the module table. This is the property that allows a tensor clause to
introduce \texttt{i} without changing the meaning of a neighboring clause.

The indexed-clause case is especially important. Output indices and
where-bound variables are allocated in an index scope before the body is
resolved, so every rectangular access and reduction body can point at the
precise DefId it uses. Later passes never need to inspect textual names to
recover that relationship. If a diagnostic still wants to print the spelling
\texttt{i}, it can do so as presentation; the compiler's semantic decisions are
already made from DefIds. That split between user-facing names and compiler
identity is one of the main reasons the rest of the implementation can remain
modular.

The ordering of the traversal is also part of the algorithm's correctness.
Hoisting items before resolving bodies means that a call can refer to a
function that appears later in the file, while local bindings remain
source-ordered so a block cannot accidentally use a value before it is
introduced. The resolver therefore uses two different temporal models: an
item-level model for modules and functions, and a statement-level model for
locals. This split is small in code, but it is what lets \einlang{} support both
ordinary function libraries and expression-like blocks without a second
resolution mechanism.

The failure modes are designed to remain source-facing. An unresolved
identifier is reported at the use site, not at a later backend lookup; an
attempt to bind an index in the wrong scope is rejected while the compiler still
knows which clause introduced it; and an imported name conflict is reported
against the relevant \texttt{use} declaration. The later compiler can then
treat missing DefIds as internal bugs rather than user errors. This is an
important thesis point: the resolver turns textual ambiguity into either a
stable identity or an early diagnostic, and every smart pass depends on that
binary outcome.

\textit{Concrete example.} In a program containing \texttt{fn norm(x) \{ let
x = x * x; x \}} and a separate indexed binding \texttt{let y[i] = x[i];}, the
printed name \texttt{x} names three different things: the function parameter,
the local square, and the top-level tensor read by \texttt{y}. Resolution gives
each binding its own DefId before type or shape analysis runs. The local
\texttt{x} shadows the parameter only inside the function body, while the
\texttt{x[i]} read in the indexed binding still points to the top-level tensor.
This is the concrete reason later passes do not carry name-disambiguation code:
they see only DefIds and scopes that have already been checked.

\begin{table}[htbp]
\centering
\caption{Representative DefId ownership rules.}
\label{tab:defid-rules}
\renewcommand{\arraystretch}{1.08}
\begin{tabular}{@{}L{0.26\textwidth}L{0.62\textwidth}@{}}
\toprule
Definition kind & Resolution behavior \\
\midrule
Builtins & Allocated in a fixed builtin crate so runtime and compiler can agree on names such as \texttt{print}, \texttt{shape}, and \texttt{\_\_basis\_tensor}. \\
Functions and constants & Registered in the module symbol table and resolver registry. They can be imported, re-exported, and looked up by module path. \\
Parameters and locals & Allocated as local DefIds and attached to definitions and use sites, but not inserted into the global symbol table. \\
Index variables & Allocated as local DefIds so repeated spellings such as \texttt{i} in different indexed bindings remain distinct. \\
Generated indices & Allocated by the resolver when rest-pattern or lowering passes need compiler-created loop variables. \\
\bottomrule
\end{tabular}
\end{table}

\section{Module Loading}

\einlang{} modules are file based. The module system under
\pathcode{src/einlang/analysis/module_system} resolves module paths, loads
source, parses module ASTs, detects circular imports, processes submodule
declarations, and handles exported symbols. The path resolver follows
Rust-inspired rules:

\begin{itemize}
\item \texttt{std::math} resolves into the bundled standard-library tree;
\item \texttt{crate::x}, \texttt{self::x}, and \texttt{super::x} are handled as
structured module paths;
\item project-local modules resolve relative to the crate root;
\item \texttt{python::...} is a virtual path for Python module integration.
\end{itemize}

Module loading is performance-sensitive because standard-library modules are
used repeatedly. The loader caches file contents and parsed ASTs in process, and
can optionally use a disk parse cache when \texttt{EINLANG\_CACHE\_DIR} is set.
Callers receive fresh AST copies because later compiler passes mutate node
metadata.

The loader also supports source overlays. Tests and embedding callers can supply
in-memory module sources keyed by module path, avoiding filesystem reads while
still exercising the same module and name-resolution code as normal programs.

\subsection{Module Loading Algorithm}

Module loading is deliberately written as a cache-and-copy algorithm. The
cached object is parsed syntax; the returned object is a fresh copy that later
passes may annotate or mutate.

\begin{lstlisting}[style=pythoncode,caption={Source overlay and module-cache resolution.}]
load_module(path, requester):
    normalized = path_resolver.resolve(path, requester)
    if normalized in active_import_stack:
        report circular import

    if normalized in source_overlays:
        source = source_overlays[normalized]
        cache_key = ("overlay", normalized, hash(source))
    else:
        file_path = path_resolver.to_file(normalized)
        source = read_file_cached(file_path)
        cache_key = ("file", file_path, mtime_or_content_key)

    if cache_key not in parsed_ast_cache:
        parsed_ast_cache[cache_key] = parse_source(source)

    ast = deep_copy(parsed_ast_cache[cache_key])
    active_import_stack.push(normalized)
    process_submodules_and_use_statements(ast)
    active_import_stack.pop()
    return ast
\end{lstlisting}

The overlay branch is important for tests and embedding: it makes synthetic
modules go through the same import, symbol-linking, and resolution machinery as
files on disk. The deep-copy step is equally important because downstream IR
and AST passes attach DefIds, types, shapes, ranges, and lowering facts.

The cache stores parsed syntax rather than analyzed syntax. That choice avoids
sharing mutable compiler annotations between compilations, which would be a
subtle source of stale DefIds, stale shape facts, or stale diagnostics. A cache
hit therefore saves parsing work but still returns a tree that the caller can
own. This is a conservative design: it gives repeated stdlib imports and tests
the performance benefit of caching without turning later passes into
copy-on-write clients of a global AST.

The active import stack is the other non-trivial part of the algorithm. Module
loading is recursive because a module may declare submodules or import other
paths while it is being processed. Recording the normalized path before
descending lets the loader report circular imports as source errors rather than
recursing indefinitely. The same normalized path also lets source overlays and
filesystem modules share the symbol-linking path, which keeps test fixtures and
real projects behaviorally aligned.

The cache key design prevents a subtle class of bugs. Overlay modules are keyed
by content hash because two tests may reuse the same logical module path with
different in-memory source. Filesystem modules are keyed by a file-derived
freshness value, so editing a standard-library file can invalidate the cached
parse tree. In both cases the cached artifact is syntax only. DefIds, types,
shapes, and module exports are deliberately recomputed because they depend on
the compilation context that requested the module.

The algorithm also makes module loading compositional. A module can import a
submodule, that submodule can import a sibling, and all of those paths collapse
to normalized module identities before caching and cycle detection run. This is
why the circular-import stack stores normalized paths rather than raw source
strings. Without normalization, \texttt{self::x}, \texttt{crate::m::x}, and a
relative filesystem spelling could look different while naming the same module.
The loader's job is to collapse those spellings before the resolver assigns
semantic identities to exported items.

\textit{Concrete example.} A source file that says \texttt{use
std::math::basic::sqrt;} first resolves the path to the bundled stdlib module,
then parses or reuses the cached AST for that module, then returns a fresh copy
for this compilation. A test may override the same logical path with a source
overlay defining a tiny fake \texttt{sqrt}; the cache key includes the overlay
contents, so that fake module cannot poison later filesystem loads. In both
cases the resolver receives an ordinary module AST and assigns DefIds to the
exported function in the requesting compilation context.

\section{Imports, Aliases, and Re-exports}

The symbol linker resolves \texttt{use} statements into function mappings and
module aliases. It supports direct imports, wildcard imports, module aliases,
and public re-exports. For example, \texttt{use std::math::basic::sqrt;} maps
the visible name \texttt{sqrt} to the module that exported it, while
\texttt{use std::math as m;} creates an alias used by later module-access
expressions.

This module machinery is not decoration around the core language. The standard
library, examples, and tests depend on it, and the implementation treats stdlib
functions as ordinary \einlang{} code rather than as compiler magic whenever
possible.
